\documentclass[12pt,a4paper]{article}
\usepackage[utf8]{inputenc}
\usepackage[indonesian]{babel}
\usepackage{amsmath}
\usepackage{amsfonts}
\usepackage{amssymb}
\usepackage{geometry}

\geometry{margin=2.5cm}

\title{Penyelesaian Volume Benda Putar dengan Dalil Guldin I}
\author{}
\date{}

\begin{document}

\maketitle

\section{Soal}
Tentukan volume benda putar yang terbentuk dari bidang datar yang dibatasi oleh kurva $y = 6 - 3x - x^2$ dan garis $x + y - 3 = 0$, jika diputar pada garis $y = 3 - x$ menggunakan Dalil Guldin I.

\section{Penyelesaian}

Dalil Guldin I menyatakan bahwa volume ($V$) benda putar adalah hasil kali antara luas daerah ($A$) dengan jarak yang ditempuh oleh titik berat daerah tersebut ($2\pi d$):
\begin{equation}
  V = 2\pi \cdot d \cdot A
\end{equation}

\subsection{Titik Potong dan Luas Daerah ($A$)}
Persamaan kurva: $y_1 = 6 - 3x - x^2$ \\
Persamaan garis: $y_2 = 3 - x$

Titik potong:
\begin{align*}
  6 - 3x - x^2 &= 3 - x \\
  -x^2 - 2x + 3 &= 0 \\
  x^2 + 2x - 3 &= 0 \\
  (x + 3)(x - 1) &= 0
\end{align*}
Batas integrasi adalah $x = -3$ dan $x = 1$.

Luas daerah $A$:
\begin{align*}
  A &= \int_{-3}^{1} [(6 - 3x - x^2) - (3 - x)] \, dx \\
  A &= \int_{-3}^{1} (-x^2 - 2x + 3) \, dx \\
  A &= \left[ -\frac{1}{3}x^3 - x^2 + 3x \right]_{-3}^{1} \\
  A &= \left( -\frac{1}{3} - 1 + 3 \right) - (9 - 9 - 9) = \frac{5}{3} + 9 = \frac{32}{3}
\end{align*}

\subsection{Koordinat Titik Berat $(\bar{x}, \bar{y})$}

\subsubsection{Menghitung $\bar{x}$}
\begin{align*}
  \bar{x} &= \frac{1}{A} \int_{-3}^{1} x(-x^2 - 2x + 3) \, dx \\
  \bar{x} &= \frac{3}{32} \int_{-3}^{1} (-x^3 - 2x^2 + 3x) \, dx \\
  \bar{x} &= \frac{3}{32} \left[ -\frac{1}{4}x^4 - \frac{2}{3}x^3 + \frac{3}{2}x^2 \right]_{-3}^{1} \\
  \bar{x} &= \frac{3}{32} \left[ \left( -\frac{1}{4} - \frac{2}{3} + \frac{3}{2} \right) - \left( -\frac{81}{4} + 18 + \frac{27}{2} \right) \right] = -1
\end{align*}

\subsubsection{Menghitung $\bar{y}$}
Menggunakan rumus $\bar{y} = \frac{1}{2A} \int (y_1^2 - y_2^2) \, dx$:
\begin{align*}
  y_1^2 - y_2^2 &= (6 - 3x - x^2)^2 - (3 - x)^2 \\
  &= (x^4 + 6x^3 - 3x^2 - 36x + 36) - (x^2 - 6x + 9) \\
  &= x^4 + 6x^3 - 4x^2 - 30x + 27
\end{align*}
\begin{align*}
  \bar{y} &= \frac{1}{2(32/3)} \int_{-3}^{1} (x^4 + 6x^3 - 4x^2 - 30x + 27) \, dx \\
  \bar{y} &= \frac{3}{64} \left[ \frac{1}{5}x^5 + \frac{3}{2}x^4 - \frac{4}{3}x^3 - 15x^2 + 27x \right]_{-3}^{1} \\
  \bar{y} &= \frac{3}{64} \left( \frac{1792}{15} \right) = \frac{28}{5} = 5,6
\end{align*}
Titik berat daerah adalah $Z(-1; 5,6)$.

\subsection{Jarak Titik Berat ke Sumbu Putar ($d$)}
Sumbu putar: $x + y - 3 = 0$.
\begin{align*}
  d &= \frac{|ax_1 + by_1 + c|}{\sqrt{a^2 + b^2}} \\
  d &= \frac{|1(-1) + 1(5,6) - 3|}{\sqrt{1^2 + 1^2}} = \frac{1,6}{\sqrt{2}}
\end{align*}

\subsection{Volume Benda Putar ($V$)}
\begin{align*}
  V &= 2\pi \cdot d \cdot A \\
  V &= 2\pi \cdot \frac{1,6}{\sqrt{2}} \cdot \frac{32}{3} \\
  V &= \frac{102,4\pi}{3\sqrt{2}} = \frac{51,2\sqrt{2}\pi}{3} \approx 24,13\pi
\end{align*}

\section{Kesimpulan}
Volume benda putar tersebut adalah $\mathbf{24,13\pi}$ satuan volume.

\end{document}